% Generated draft: requires mathematical and source review.
\documentclass[11pt]{article}
\usepackage[margin=25mm]{geometry}
\usepackage{fontspec}
\setmainfont{Latin Modern Roman}
\usepackage{amsmath,amssymb,mathtools,unicode-math}
\setmathfont{Latin Modern Math}
\usepackage{enumitem,needspace,xurl,hyperref}
\setlength{\parindent}{0pt}
\setlength{\parskip}{5pt}
\setlength{\emergencystretch}{3em}
\newcommand{\N}{\mathbb{N}}
\newcommand{\Z}{\mathbb{Z}}
\newcommand{\Q}{\mathbb{Q}}
\newcommand{\R}{\mathbb{R}}
\newcommand{\C}{\mathbb{C}}
\newcommand{\F}{\mathbb{F}}
\newcommand{\A}{\mathbb{A}}
\newcommand{\PP}{\mathbb{P}}
\newcommand{\E}{\mathbb{E}}
\newcommand{\eps}{\varepsilon}
\newcommand{\dd}{\,\mathrm{d}}
\begin{document}
\section*{UnsolvedMath ID: 30001228}
\textbf{Source problem code:} catalogue entry
\subsection{Definitions and mathematical statement}
Let \(d\in\N\) with \(d\geq 3\). Let \(H\) be a Hilbert space and let \(\mathcal B(H)\) denote its bounded linear operators. A commuting \(d\)-tuple of contractions on \(H\) is a tuple \(T=(T_1,\ldots,T_d)\in\mathcal B(H)^d\) satisfying
\[
\|T_i\|\leq 1\quad(1\leq i\leq d),\qquad T_iT_j=T_jT_i\quad(1\leq i,j\leq d).
\]
Here \(d\) is the number of operators, not the dimension of \(H\). The supplied record states no restriction on the dimension or separability of \(H\), and does not explicitly specify its scalar field. All spaces in an extension are understood to have the same scalar field; the intended choice of real or complex scalars remains to be confirmed.

Using the usual invariant-subspace meaning of extension, an extension of \(T\) within this family is a commuting \(d\)-tuple of contractions \(R=(R_1,\ldots,R_d)\) on a Hilbert space \(K\) containing \(H\) as a closed subspace, such that
\[
R_iH\subseteq H,\qquad R_i|_H=T_i\quad(1\leq i\leq d).
\]
A closed subspace \(H\subseteq K\) reduces \(R\) if it is invariant under every \(R_i\) and every adjoint \(R_i^*\). The tuple \(T\) is called extremal if, for every such \(K\) and every such extension \(R\), the subspace \(H\) reduces \(R\). Equivalently, every extension has, relative to \(K=H\oplus H^\perp\), the form
\[
R_i=T_i\oplus S_i\quad(1\leq i\leq d),
\]
where \(S=(S_1,\ldots,S_d)\) is a commuting tuple of contractions on \(H^\perp\).

\textbf{Problem.} For every integer \(d\geq 3\), characterize the commuting \(d\)-tuples of contractions that are extremal in this sense. The scalar-field convention and the interpretation of extension above require confirmation against the source report. The supplied open status is a dated curation label, not an independently verified assertion of present status.
\subsection{Short English statement}
For each integer \(d>2\), characterize commuting tuples of \(d\) contractions for which every extension to a larger Hilbert space splits as a direct sum with the original tuple. Thus no extension may couple the original space to its orthogonal complement. The supplied record leaves the scalar field unspecified; the definition uses the usual invariant-subspace interpretation of extension.
\subsection{Sources}
\begin{itemize}
\item Supplied local UnsolvedMath registry record, permanent ID \(30001228\), titled \emph{Extremal Tuples of Commuting Contractions}; collection: Oberwolfach Reports Open Problems; exact source problem code: \texttt{catalogue entry}. The recorded proposed year is 2009; no individual proposer is supplied.
\item \emph{Hilbert Modules and Complex Geometry} (2009), Oberwolfach Reports, pp.\ 1047--1100, DOI \texttt{10.4171/owr/2009/20}. Supplied report reference and literature-source link, labeled ``original report'': \url{https://doi.org/10.4171/owr/2009/20}. The supplied background describes Oberwolfach Reports as documenting mathematical workshops and their open-problem discussions.
\item Supplied literature assessment, checked on \texttt{2026-08-21}: status \emph{open}. The assessment says that the report explicitly states that extremals for commuting contraction tuples with more than two components were unknown. This is supplied dated curation derived from the OWR record, not independently verified literature research; time-sensitive status and cited claims require independent verification.
\item The registry labels the statement \texttt{corrected\_verified}, with the supplied explanation that its semantic, typographic, and LaTeX cleanup preserved meaning and added no unsupported source-dependent content. This label is recorded as supplied metadata, not as independent mathematical verification of this draft.
\end{itemize}
\end{document}
